
\documentclass{article}
\usepackage{colortbl}
\usepackage{makecell}
\usepackage{multirow}
\usepackage{supertabular}

\begin{document}

\newcounter{utterance}

\centering \large Interaction Transcript for game `cladder', experiment `full\_v1.5\_default', episode 6572 with refineonly{-}llama{-}3.1{-}8b{-}instruct.
\vspace{24pt}

{ \footnotesize  \setcounter{utterance}{1}
\setlength{\tabcolsep}{0pt}
\begin{supertabular}{c@{$\;$}|p{.15\linewidth}@{}p{.15\linewidth}p{.15\linewidth}p{.15\linewidth}p{.15\linewidth}p{.15\linewidth}}
    \# & \multicolumn{2}{c}{Player} && \multicolumn{2}{c}{Game Master} \\
    \hline

    \theutterance \stepcounter{utterance}  
    & & & \multicolumn{4}{p{0.6\linewidth}}{
        \cellcolor[rgb]{0.9,0.9,0.9}{
            \makecell[{{p{\linewidth}}}]{
                \texttt{\tiny{[P1$\langle$GM]}}
                \texttt{You are an expert in causal inference. The following question is not a typical commonsense query, but rather a meticulously designed question created by a professor specializing in causal inference, intended to assess the students' mastery of the course content.} \\
\texttt{Imagine a self{-}contained, hypothetical world with only the following conditions, and without any unmentioned factors or causal relationships: Kwox has a direct effect on kwoz. Swoq has a direct effect on kwoz. For those who are not kwox, the probability of kwoz is 82\%. For those who are kwox, the probability of kwoz is 32\%. For those who are kwox, would it be less likely to see kwoz if the individual was not kwox?} \\
\texttt{Start your answer with 'Yes' or 'No', followed by additional reasoning or evidence to support your explanation.} \\
            }
        }
    }
     \\ \\

    \theutterance \stepcounter{utterance}  
    & \multicolumn{4}{p{0.6\linewidth}}{
        \cellcolor[rgb]{0.561,0.737,0.561}{
            \makecell[{{p{\linewidth}}}]{
                \texttt{\tiny{[P1$\rangle$GM]}}
                \texttt{No.} \\
\\ 
\texttt{The counterfactual question is asking whether it would be less likely to see kwoz if the individual were not Kwox, given that they are Kwox. This is equivalent to asking whether the probability of kwoz would be higher if the individual were not Kwox, given that they are Kwox.} \\
\\ 
\texttt{Using the law of total probability, we can calculate the probability of kwoz given that the individual is Kwox and not Kwox. Let P(Kwox) be the probability of being Kwox, and P(kwoz|Kwox) be the probability of kwoz given that the individual is Kwox. Then, the probability of kwox and kwoz is P(Kwox) * P(kwoz|Kwox) = 0.5 * 0.32 = 0.16.} \\
\\ 
\texttt{However, we are given that for those who are not Kwox, the probability of kwoz is 82\%. This means that the probability of kwoz given that the individual is not Kwox is P(kwoz|not Kwox) = 0.82.} \\
\\ 
\texttt{To calculate the probability of kwoz given that the individual is Kwox and not Kwox, we can use the law of total probability:} \\
\\ 
\texttt{P(kwoz|Kwox and not Kwox) = P(kwoz|Kwox) * P(not Kwox|Kwox) + P(kwoz|not Kwox) * P(not Kwox|Kwox)} \\
\\ 
\texttt{Since P(not Kwox|Kwox) = 1 {-} P(Kwox|Kwox) = 1 {-} 1 = 0 (which is not possible), we cannot calculate the probability of kwoz given that the individual is Kwox and not Kwox using this method.} \\
\\ 
\texttt{However, we can use the concept of counterfactuals to answer the question. The counterfactual question is asking whether it would be less likely to see kwoz if the individual were not Kwox, given that they are Kwox. This is equivalent to asking whether the probability of kwoz would be higher if the individual were not Kwox, given that they are Kwox.} \\
\\ 
\texttt{Since we are given that for those who are not Kwox, the probability of kwoz is 82\%, it is likely that the probability of kwoz would be higher if the individual were not Kwox, given that they are Kwox.} \\
            }
        }
    }
    & & \\ \\

    \theutterance \stepcounter{utterance}  
    & & & \multicolumn{2}{p{0.3\linewidth}}{
        \cellcolor[rgb]{0.9,0.9,0.9}{
            \makecell[{{p{\linewidth}}}]{
                \texttt{\tiny{[GM$|$GM]}}
                \texttt{No} \\
            }
        }
    }
    & & \\ \\

    \theutterance \stepcounter{utterance}  
    & & & \multicolumn{2}{p{0.3\linewidth}}{
        \cellcolor[rgb]{0.9,0.9,0.9}{
            \makecell[{{p{\linewidth}}}]{
                \texttt{\tiny{[GM$|$GM]}}
                \texttt{no} \\
            }
        }
    }
    & & \\ \\

    \theutterance \stepcounter{utterance}  
    & & & \multicolumn{2}{p{0.3\linewidth}}{
        \cellcolor[rgb]{0.9,0.9,0.9}{
            \makecell[{{p{\linewidth}}}]{
                \texttt{\tiny{[GM$|$GM]}}
                \texttt{game\_result = WIN} \\
            }
        }
    }
    & & \\ \\

\end{supertabular}
}

\end{document}
